% =====================================================================
%  facilitator_manual.tex
%
%  Facilitator's Manual for the Security Cards Solutions Dimension.
%  Standalone document -- does NOT share a preamble with main.tex.
%  Build: pdflatex facilitator_manual (x2 for the table of contents)
%
%  This is a working document for classroom use, not a submission.
%  Edit freely; the section most worth localizing to your course is
%  Section 7 (intervention) and Appendix A (session record).
% =====================================================================

\documentclass[11pt,letterpaper]{article}

\usepackage[letterpaper,margin=1in,headheight=14pt]{geometry}
\usepackage{mathptmx}
\usepackage[T1]{fontenc}
\usepackage[utf8]{inputenc}
\usepackage{microtype}
\usepackage{amssymb}
\usepackage{graphicx}
\usepackage{booktabs}
\usepackage{tabularx}
\usepackage{longtable}
\usepackage{array}
\usepackage{xcolor}
\usepackage{enumitem}
\usepackage{caption}
\usepackage{fancyhdr}
\usepackage[most]{tcolorbox}
\usepackage{titlesec}
\usepackage[hidelinks]{hyperref}

\definecolor{deep}{HTML}{1F3A5F}
\definecolor{warn}{HTML}{9C3B2E}
\definecolor{calm}{HTML}{2F6B4F}
\definecolor{soft}{HTML}{F2F4F7}

\titleformat{\section}{\normalfont\Large\bfseries\color{deep}}{\thesection.}{0.5em}{}
\titleformat{\subsection}{\normalfont\large\bfseries}{\thesubsection.}{0.5em}{}
\titlespacing*{\section}{0pt}{2.4ex}{1.1ex}

\pagestyle{fancy}
\fancyhf{}
\lhead{\small\color{deep}Facilitator's Manual --- Solutions Dimension}
\rhead{\small\thepage}
\renewcommand{\headrulewidth}{0.4pt}

\setlist{itemsep=2pt,parsep=0pt,topsep=3pt}
\renewcommand{\arraystretch}{1.25}
\newcolumntype{R}{>{\raggedright\arraybackslash}X}
\captionsetup{font=small,labelfont=bf}

\newtcolorbox{intervene}[1]{colback=warn!6,colframe=warn,
  fonttitle=\bfseries,title=#1,boxrule=0.8pt,arc=2pt,left=6pt,right=6pt,
  top=4pt,bottom=4pt}
\newtcolorbox{hold}[1]{colback=calm!6,colframe=calm,
  fonttitle=\bfseries,title=#1,boxrule=0.8pt,arc=2pt,left=6pt,right=6pt,
  top=4pt,bottom=4pt}
\newtcolorbox{note}[1]{colback=soft,colframe=deep!45,
  fonttitle=\bfseries,title=#1,boxrule=0.6pt,arc=2pt,left=6pt,right=6pt,
  top=4pt,bottom=4pt}

\newcommand{\say}[1]{\textit{``#1''}}

\begin{document}

\begin{titlepage}
\centering
\vspace*{2.2cm}
{\color{deep}\rule{\textwidth}{2pt}}\\[0.8em]
{\Huge\bfseries\color{deep} Facilitator's Manual}\\[0.6em]
{\LARGE The Solutions Dimension}\\[0.4em]
{\large An extension to the University of Washington Security Cards}\\[0.5em]
{\color{deep}\rule{\textwidth}{2pt}}\\[2.2cm]

{\large 28 Solutions cards $\cdot$ 4 optional inject cards $\cdot$ 6 gameplay variants}\\[0.4em]
{\large Mapped to CIS Critical Security Controls v8.1}\\[3cm]

\begin{minipage}{0.82\textwidth}
\begin{note}{Who this manual is for}
Anyone running the activity: course instructors, teaching assistants,
clinic staff, or students facilitating for their peers. \textbf{No security
expertise is required to facilitate.} Everything you need to adjudicate is in
Sections 6 and 7. If you read only two sections before your first session, read
Section 3 (Quick Reference) and Section 7 (When to Intervene).
\end{note}
\end{minipage}

\vfill
{\small Version 2.0 \quad$\cdot$\quad Deck released under an open license}
\end{titlepage}

\tableofcontents
\newpage

% =====================================================================
\section{What This Game Is, and Is Not}

The Security Cards were built to help groups think broadly about security
threats. They stop at threat discovery on purpose: there is no scoring, no
budget, and no defined end state. The Solutions dimension adds the other half of
the professional task---choosing what to actually do about a threat, with
limited money.

\textbf{What the activity teaches.} Prioritization under constraint. Students
practice moving from a stated threat to a proportionate, affordable control, and
justifying that choice out loud to someone who may disagree.

\textbf{What it does not teach.} It is not a simulation, a risk assessment
methodology, or a substitute for technical instruction. A student who plays well
is not thereby competent to advise a client. The activity is an on-ramp: it
gives learners a vocabulary, a framework identifier they can look up afterward,
and practice at one specific cognitive move.

\begin{note}{The single most important thing to understand}
\textbf{The disagreement is the lesson.} Students learning the card definitions
is a side effect. The activity works when a team argues about which of two
defensible options is better and has to articulate why. A session where everyone
agreed quickly is a session that did not work, and Section 7 tells you what to
do about it.
\end{note}

% =====================================================================
\section{Before the Session}

\subsection{Materials}

\begin{itemize}
\item \textbf{Original Security Cards deck} (42 cards, four dimensions), one per
team. Teams may share if you are short.
\item \textbf{Solutions deck}, one per team: 21 control cards, 7 modifier cards.
\item \textbf{Risk matrix}, one per team. A printed $3\times3$ grid, or drawn on
paper. Large enough to lay cards inside the cells.
\item \textbf{Budget tokens}, 6 per team. Poker chips, paper squares, or coins.
Physical tokens work substantially better than tallying on paper, because
spending becomes visible and irreversible.
\item \textbf{Session record sheet} (Appendix A), one per team.
\item A timer everyone can see.
\end{itemize}

\subsection{Printing and Setup}

Print Solutions cards double-sided on card stock. If printing in grayscale,
confirm that Solutions cards remain distinguishable from Human Impact cards ---
both use a blue family, and the icon frame shape is the backup differentiator.

Sort each Solutions deck into three piles before the session: Foundational
(cards 1--13), Extended (14--18), Advanced (19--21), with modifiers (22--28)
separate. You will deal from these piles rather than shuffling everything
together.

\subsection{Room and Group Size}

Teams of \textbf{two to three}. Four is the maximum before someone stops
talking. Seat teams so you can circulate to every table without walking behind
anyone's chair.

For class size, divide by three and round up. The web version defaults to six
teams and can be set to 8, 10, 12 or 16 from the facilitator panel, which covers
a room of about sixty-four. Reducing the count later hides the higher-numbered
teams rather than deleting them, so nothing is lost if you guess high.

\textbf{One device drives per team.} The board is shared, so if two people on the
same team click at the same moment one set of changes overwrites the other.
Everyone may join and watch --- the app shows how many devices are on a team, and
warns anyone who is not the driver --- but agree who is typing before Phase 2
starts. On paper this problem does not exist, which is one reason the physical
deck still generates more discussion.

Multiple teams playing simultaneously in one room works well and is the
recommended arrangement, because the cross-team comparison at the end is where
much of the learning lands. Expect it to be loud.

\subsection{Choosing a Scenario}

Do not improvise one. An invented scenario is the most common reason a session
goes flat, and it also destroys cross-team comparison, since two teams arguing
about different situations are not disagreeing about anything.
Appendix~\ref{app:scenarios} has ten ready to read aloud, each with an expert key
and its own debrief questions. Scenarios 1, 2, 4 and 8 suit a first session.

\subsection{Timing}

\begin{center}
\small
\begin{tabular}{@{}l l l@{}}
\toprule
\textbf{Segment} & \textbf{Time} & \textbf{Note} \\
\midrule
Setup and rules explanation & 5--7 min & Do not over-explain; see below \\
Phase 1: threat and placement & 5--7 min & \\
Phase 2: budgeted defense & 8--10 min & The longest phase \\
Phase 3: residual risk & 4--5 min & \\
Cross-team comparison & 5--8 min & Do not skip this \\
Debrief & 5--10 min & Section 10 \\
\midrule
\textbf{Total} & \textbf{35--45 min} & Fits a 50-minute period \\
\bottomrule
\end{tabular}
\end{center}

\begin{note}{On explaining the rules}
Explain Phase 1 only, then start. Explain Phase 2 when teams reach it, and
Phase 3 when they reach that. Front-loading all three phases produces confused
teams and wastes five minutes. The game is designed to be learned in motion.
\end{note}

% =====================================================================
\section{Quick Reference}

Keep this page in front of you during the session.

\begin{center}
\small
\begin{tabularx}{\textwidth}{@{}l X@{}}
\toprule
\textbf{Deal} & 8 control cards + 3 modifier cards per team \\
\textbf{Budget} & 6 tokens per team (see Section 8 for other tiers) \\
\textbf{Costs} & Foundational = 1 \quad Extended = 2 \quad Advanced = 3 \quad Modifiers = 0 \\
\textbf{Modifiers} & At most 2 played per round, each requires a stated justification \\
\textbf{Matrix} & $3\times3$: likelihood (Not Likely / Likely / Very Likely) $\times$ severity (Not Severe / Severe / Very Severe) \\
\midrule
\textbf{Phase 1} & Draw 1 card from each of the 4 original dimensions. Build a scenario. Place on matrix. \emph{Severity must be justified by naming the Human Impact card.} \\
\textbf{Phase 2} & Deal 8 Solutions + 3 modifiers. Spend up to 6. For each purchase, say which safeguard addresses which part of the threat. \\
\textbf{Phase 3} & Re-place the same threat given what was bought. State what still gets through. \\
\midrule
\textbf{Your job} & Time, budget arithmetic, modifier adjudication, and the interventions in Section 7. Nothing else. \\
\bottomrule
\end{tabularx}
\end{center}

\begin{center}
\small
\begin{tabularx}{\textwidth}{@{}c l c X@{}}
\toprule
\textbf{\#} & \textbf{Card} & \textbf{Cost} & \textbf{CIS Control} \\
\midrule
1 & Software Updates & 1 & 7 Continuous Vulnerability Mgmt \\
2 & Social Engineering Awareness & 1 & 14 Security Awareness \\
3 & Deepfake \& AI Fraud Awareness & 1 & 14 Security Awareness \\
4 & Identity \& Access Management & 1 & 6 Access Control Mgmt \\
5 & Password \& Authentication & 1 & 5, 6 Account / Access Control \\
6 & Backup \& Recovery & 1 & 11 Data Recovery \\
7 & Account Lifecycle Management & 1 & 5 Account Management \\
8 & Physical Security & 1 & --- (supports 1, 3) \\
9 & Asset Inventory & 1 & 1 Enterprise Asset Inventory \\
10 & Software Inventory & 1 & 2 Software Asset Inventory \\
11 & Secure Configuration & 1 & 4 Secure Configuration \\
12 & Email \& Browser Protections & 1 & 9 Email \& Browser \\
13 & Incident Response Planning & 1 & 17 Incident Response Mgmt \\
14 & Encryption & 2 & 3 Data Protection \\
15 & Audit Logging \& SIEM & 2 & 8, 13 Log Mgmt / Monitoring \\
16 & Endpoint Detection \& Response & 2 & 10, 13 Malware / Monitoring \\
17 & Third-Party \& Contractual & 2 & 15 Service Provider Mgmt \\
18 & Network Segmentation & 2 & 12, 13 Network Infrastructure \\
19 & Penetration Testing & 3 & 18 Penetration Testing \\
20 & Managed Detection \& Response & 3 & 8, 13, 17 \\
21 & Cyber Insurance \& Risk Transfer & 3 & --- (risk transfer) \\
\midrule
22 & Defense in Depth & 0 & Modifier \\
23 & Fail Safe Defaults & 0 & Modifier \\
24 & Least Privilege & 0 & Modifier (requires 4 or 7) \\
25 & Human Factors & 0 & Modifier \\
26 & Risk Acceptance & 0 & Modifier (extended) \\
27 & Compensating Control & 0 & Modifier (extended) \\
28 & Deferred Investment & 0 & Modifier (campaign play) \\
\bottomrule
\end{tabularx}
\end{center}

% =====================================================================
\section{The Core Loop (Variant V6)}

This is the recommended variant and the one everything else in this manual
assumes.

\subsection{Phase 1 --- Threat and Placement (5--7 minutes)}

Each team draws one card from each of the four original dimensions: Human
Impact, Adversary's Motivations, Adversary's Resources, Adversary's Methods.
They construct a short scenario in which all four appear, then place that threat
on the risk matrix.

\textbf{Say:} \say{Build a story where all four of these cards are true at once.
It does not need to be realistic in every detail, but every card has to do some
work in it.}

\textbf{The severity rule.} Teams must justify their severity placement by
pointing at the Human Impact card and saying who specifically is harmed and how
badly. This matters more than it looks. Novices rate severity by how
sophisticated an attack \emph{sounds}, so a technically impressive attack that
harms nobody gets marked Very Severe while a simple one that exposes a client's
home address gets marked Not Severe. Requiring them to point at the Human Impact
card corrects this without a lecture.

\textbf{Record} the coordinates on the session sheet. You will need them in
Phase 3.

\subsection{Phase 2 --- Budgeted Defense (8--10 minutes)}

Deal each team eight Solutions cards and three modifiers, and give them six
tokens.

Deal from the sorted piles rather than at random. A workable default is six
Foundational, one Extended, and one Advanced. If a team's threat obviously calls
for a particular card, include it --- the deal is a teaching instrument, not a
lottery.

Teams purchase a portfolio. For each purchase they must say which safeguard
addresses which element of the threat.

\textbf{Say:} \say{Six tokens. Anything you do not spend is gone. For each card
you buy, tell me which part of your scenario it touches.}

\textbf{Enforce the arithmetic immediately.} If a team overspends, correct it the
moment you notice. The constraint is the entire mechanic, and a team that
overspends unnoticed has not played the game.

\subsection{Phase 3 --- Residual Risk (4--5 minutes)}

The team re-places the same threat on the matrix given what it bought, and
states what would still get through.

\textbf{Say:} \say{Same threat, same board. Where does it sit now that you have
bought these? And tell me what still gets through, because something does.}

This is where the central lesson lands: controls reduce risk, they do not
eliminate it. A team that moves its threat to Not Likely / Not Severe is almost
always wrong, and the right response is a question rather than a correction ---
see Section 7.

\subsection{Cross-Team Comparison (5--8 minutes)}

Do not skip this. Have two teams that faced different threats state what they
bought and why. The interesting question is never who was right; it is why two
groups of reasonable people bought different things.

% =====================================================================
\section{The Other Five Variants}

\begin{center}
\small
\begin{tabularx}{\textwidth}{@{}l X l@{}}
\toprule
\textbf{Variant} & \textbf{What it does and when to use it} & \textbf{Time} \\
\midrule
V1 Scenario Generation & Four dimension cards, build an incident narrative. Use in week one, before Solutions cards are introduced at all. Lowest barrier of any variant. & 8--10 min \\
V2 Risk Matrix Placement & Rank the seven Methods cards on the matrix for a named organization. Use to introduce likelihood and severity as separate axes. Pre-populate two or three cells to lower the entry cost. & 10--15 min \\
V3 Solution Ranking & Deal 8 Solutions, rank the top three for a stated scenario. No budget. Use when you have under fifteen minutes. & 7--10 min \\
V4 Control Mapping & Given an incident, identify which control failed. Use after students have seen the CIS framework once. Works well with real breach reports. & 10 min \\
V5 Budgeted Defense & Phase 2 of the core loop, standalone. Use when you want the trade-off lesson without the matrix. & 15--20 min \\
V6 Residual Risk Loop & The full three-phase loop. The recommended variant. & 20--25 min \\
\bottomrule
\end{tabularx}
\end{center}

\begin{note}{Sequencing across a course}
V1 in week one. V2 once the risk matrix has been introduced in lecture. V3 or V4
as a short mid-term activity. V6 once students have played at least one earlier
variant. Do not open with V6; it has three phases and a budget, and cold teams
stall.
\end{note}

% =====================================================================
\section{Adjudicating Modifier Cards}

Modifiers cost nothing, which is why each one carries a justification
requirement. \textbf{A modifier may be played only if the team supplies the
justification the card demands. If it cannot, the card returns to the deck.}
This is a rule, not a judgment call, and enforcing it is your job.

You do not need security expertise to adjudicate. The test in every case is
whether the team said something \emph{specific} rather than something
\emph{general}. Below, each modifier has an example of an answer to accept and
one to reject.

\begin{center}
\small
\begin{longtable}{@{}p{2.5cm} p{5.15cm} p{5.15cm}@{}}
\toprule
\textbf{Modifier} & \textbf{Accept} & \textbf{Reject} \\
\midrule
\endhead
22 Defense in Depth
& \say{Email filtering stops the message; awareness training catches it if the
filter misses. One is technical, one is human.}
& \say{We bought two things, so it is layered.} Two purchases failing the
same way are one control bought twice. \\
23 Fail Safe Defaults
& \say{If the DNS filter goes down, the site resolves normally and staff can
reach it, so it fails open. We would want it to block by default.}
& \say{It would fail safely.} No mechanism named. \\
24 Least Privilege
& \say{The volunteer coordinator sees donor payment records. She would lose
that; she needs only names and shift times.}
& \say{Everyone should have less access.} No named person, no named loss. \\
25 Human Factors
& \say{Staff will share the one MFA phone at the front desk. We would need
per-person enrollment.}
& \say{People are the weakest link.} A slogan, not a mechanism. \\
26 Risk Acceptance
& \say{We accept the risk of a physical break-in this year. The executive
director owns that decision and it is reviewed in twelve months.}
& \say{We accept the rest of the risks.} No named risk, no named owner. A risk
nobody owns has been ignored, not accepted. \\
27 Compensating Control
& \say{We cannot afford SIEM. We are substituting a monthly manual review of
admin logins. It will not catch anything in real time.}
& \say{We will do it manually instead.} No statement of what is not covered. \\
28 Deferred Investment
& \say{We are capping at four and banking two. Meanwhile we are exposed to
credential phishing for another year, and we would tell the client that.}
& \say{We will buy it next year.} No stated exposure, no duration. \\
\bottomrule
\end{longtable}
\end{center}

\begin{note}{The one-sentence test}
If you are unsure, ask: \say{Can you say that again naming a specific person,
system, or failure?} If they can, accept. If the answer stays general after one
prompt, the card returns to the deck. Do not prompt twice --- the requirement is
the point.
\end{note}

% =====================================================================
\section{When to Intervene}

This is the section most worth reading twice.

\subsection{Your Default Is Non-Intervention}

New facilitators intervene far too early. A team sitting in silence for thirty
seconds is usually thinking. A team arguing is usually learning. A team reaching
a conclusion you consider wrong, and being corrected by a peer, has learned more
than a team corrected by you.

\textbf{Your standing job is small:} keep time, enforce the budget arithmetic,
adjudicate modifiers, and prevent the four situations in Section 7.4. Everything
else is optional and most of it is counterproductive.

\subsection{The Escalation Ladder}

When you do intervene, use the lightest rung that works, and stop there.

\begin{enumerate}[label=\textbf{Rung \arabic*.}, leftmargin=4.2em]
\item \textbf{Wait.} Count to twenty before deciding a silence is a stall.
\item \textbf{Ask an open question.} \say{What is the worst thing that happens
in your scenario?} Redirects without supplying content.
\item \textbf{Point at a card.} Physically point at the Methods card and say
nothing. Astonishingly effective.
\item \textbf{State a rule.} \say{Remember you have to say which part of the
threat each purchase touches.}
\item \textbf{Give direct instruction.} Only when time is nearly gone or the
team is genuinely stuck. Supplying an answer ends the learning for that team.
\end{enumerate}

\subsection{Intervention Triggers}

\begin{center}
\small
\begin{longtable}{@{}p{3.1cm} p{2.1cm} p{7.6cm}@{}}
\toprule
\textbf{Signal} & \textbf{Threshold} & \textbf{What to do} \\
\midrule
\endhead
Total stall, nobody speaking & 60 seconds & Rung 2. \say{Read your Methods card
out loud and tell me who it hurts.} Reading aloud restarts almost every stalled
table. \\
One person doing all the talking & Two full minutes & Rung 2, addressed to a
specific other person: \say{Which card would you defend that nobody has picked
yet?} Do not name the dominant speaker or ask them to stop. \\
Agreement reached in under 90 seconds & Immediate & Rung 2. Argue the other
side: \say{Make the case for the card you just dismissed.} Fast consensus means
the trade-off was not felt. \\
Technical jargon that others do not share & First occurrence & Rung 2, directed
at the speaker: \say{Say that again for someone who has not taken a security
course.} Protects the non-majors without embarrassing anyone. \\
Buying cards without reference to the threat & Second purchase & Rung 3. Point
at the Methods card. If needed, Rung 4: \say{Which part of the scenario does
that touch?} \\
Overspending the budget & Immediate & Rung 4. Correct at once. The constraint is
the mechanic. \\
Modifier played without justification & Immediate & Rung 4. One prompt, then the
card returns to the deck (Section 6). \\
Threat moved to Not Likely / Not Severe in Phase 3 & Immediate & Rung 2.
\say{So this can no longer happen at all? What would an attacker with more time
do?} Almost always reveals the team has over-credited its controls. \\
\say{Just tell us the right answer} & After one deflection & Rung 2, then be
honest: \say{There is a better and a worse answer here, but not one right one.
Tell me your reasoning and I will tell you where it is weak.} Novices need to
hear that the ambiguity is deliberate. \\
Discussion drifting into how to perform the attack & First occurrence & Rung 4.
\say{We are on the defending side of the table.} Redirect without moralizing. \\
Team finishing five or more minutes early & Immediate & Rung 2. \say{Your budget
was cut by two tokens. What goes?} A live budget cut is the best use of spare
time. \\
Running out of time & 3 min remaining & Rung 5. Have them commit to a portfolio
even if unfinished; an incomplete decision is still assessable, an abandoned one
is not. \\
\bottomrule
\end{longtable}
\end{center}

\subsection{Four Situations That Require Immediate Intervention}

These four are different in kind from the table above. They are not pedagogical
adjustments, and the escalation ladder does not apply --- act at once.

\begin{intervene}{1. A participant discloses a real vulnerability}
Someone says their employer, their student organization, or their family
business has a specific unaddressed weakness. \textbf{Stop the disclosure at the
table.} Say: \say{Let us not work through a real system in the open --- come
find me after.} Then follow up privately and route it appropriately. Discussing
a live vulnerability in a room of twenty people creates risk for a third party
who did not consent to the conversation.
\end{intervene}

\begin{intervene}{2. A scenario lands on someone's personal experience}
The Human Impact dimension includes Physical Wellbeing, Emotional Wellbeing, and
Relationships, and the Methods dimension includes Manipulation or Coercion. A
scenario about stalking, intimate partner surveillance, doxxing, or harassment
can land hard on a participant with that history. Watch for someone going quiet,
withdrawing, or leaving.

\textbf{Offer an exit without requiring an explanation:} \say{This team can
redraw the Impact card if you want a different scenario.} Address the team, not
the individual. Never require anyone to narrate a scenario aloud. Have a
redraw available at every session as a standing, unremarkable option.
\end{intervene}

\begin{intervene}{3. A participant is being spoken over or dismissed}
Frequently a non-major being dismissed by a technical student. Intervene
immediately and structurally rather than socially: \say{Let us hear each person's
top card before anyone argues.} Restructuring the turn order fixes this without
anyone being reprimanded, and it protects the design goal that non-specialists
can play.
\end{intervene}

\begin{intervene}{4. Factual error that will propagate}
Most wrong answers should be left for peers or the debrief. Intervene at once
only when the error would be carried out of the room as a belief --- for
example, that encryption protects against phishing, that backups prevent a
breach, or that insurance removes the obligation to notify. Correct briefly and
without ceremony, then hand the discussion back.
\end{intervene}

\subsection{When Not to Intervene}

\begin{hold}{Leave these alone}
\begin{itemize}
\item \textbf{Silence under twenty seconds.} It is thinking.
\item \textbf{Disagreement, including heated disagreement.} It is the mechanism.
Only intervene if it becomes personal rather than substantive.
\item \textbf{A defensible choice you would not have made.} There are several
good answers; yours is one of them. Ask why, do not correct.
\item \textbf{A wrong answer another student is about to catch.} Wait three
seconds. Peer correction is worth more than yours.
\item \textbf{A team ignoring a card you think is obviously right.} Note it for
the debrief. Steering purchases hollows out the exercise.
\item \textbf{Slightly unrealistic scenarios.} The scenario is scaffolding, not
the object of assessment. Do not audit plausibility.
\item \textbf{A team spending everything on one expensive card.} That is a
legitimate strategy with a real consequence. Let Phase 3 teach it.
\end{itemize}
\end{hold}

% =====================================================================
\section{Campaign Play Across a Term}
\label{sec:campaign-manual}

Campaign play revisits the same client organization multiple times --- typically
weeks 3, 8, and 14. It is where the deck teaches roadmapping rather than
one-shot selection. Use it only after teams have played the core loop at least
once.

\subsection{Budget by Organization}

Budget is a property of the organization in the scenario, not of the game.

\begin{center}
\small
\begin{tabular}{@{}c l l@{}}
\toprule
\textbf{Budget} & \textbf{Organization modeled} & \textbf{CIS tier} \\
\midrule
5 & Volunteer-run, no IT staff & Below \textsc{ig1} \\
\textbf{6} & \textbf{Default; small nonprofit} & \textbf{\textsc{ig1}} \\
10 & One IT staff member & \textsc{ig2} \\
15 & Dedicated security function & \textsc{ig3} \\
\bottomrule
\end{tabular}
\end{center}

Running the same threat at two budget levels in consecutive sessions is one of
the most effective things you can do with this deck. Students discover that the
right answer changes with resources, which is the core insight of the
Implementation Group model.

\subsection{Deferred Investment and Threat Drift}

A team playing \textbf{card 28 (Deferred Investment)} caps its spend at four,
banks two, and receives ten points at the next engagement --- the \textsc{ig2}
budget. In exchange, every threat left unmitigated drifts one cell toward
\emph{Very Likely} when the scenario is revisited.

\textbf{Likelihood drifts. Severity does not.} Severity is a property of who is
harmed, which does not change because an organization delayed. Record drift on
the session sheet; you are the bookkeeper across terms.

\subsection{Realization}

Because the matrix has only three likelihood states, a deferred threat reaches
\emph{Very Likely} within two engagements. A threat that survives \emph{another}
engagement unmitigated \textbf{realizes}.

\textbf{Announce it. Do not roll dice.} The outcome is determined entirely by
what the team purchased in earlier engagements:

\begin{center}
\small
\begin{tabularx}{\textwidth}{@{}l X X@{}}
\toprule
\textbf{Question} & \textbf{Purchased} & \textbf{Not purchased} \\
\midrule
Detection (15 or 20) & Discovered within days, internally & Discovered months
later, by a third party or the press \\
Response (13) & They know who to call and in what order & Forfeit the first
purchase of the next engagement to confusion \\
Recovery (6) & Data and systems restored & The loss is permanent \\
Financial (21) & Part of the cost is transferred & The full cost lands on the
organization \\
Disclosure (17) & The provider is obliged to notify & They learn from someone
else \\
\bottomrule
\end{tabularx}
\end{center}

\textbf{The required artifact.} The team narrates the incident to the room in
the client's terms, using the Human Impact card to say who was harmed and how.
This narration is the assessable product.

\begin{note}{Why realization exists}
Twenty-one of the twenty-eight cards are preventive, and novices spend on
prevention first. A team that never experiences an incident never learns why
detection, response, and recovery were worth points. Realization supplies that
lesson retroactively without adding a card --- their earlier portfolio decides
how bad the day is.
\end{note}

\begin{intervene}{Facilitator caution on realization}
Realization is dramatic and can read as punishment. Frame it as consequence, not
failure: \say{This is what deferring looked like. Most organizations make
exactly this trade.} A team that feels punished stops taking risks in later
sessions, which defeats the purpose.
\end{intervene}

% =====================================================================
\section{Optional Inject Cards}

Four inject cards introduce adversity after purchase. \textbf{You play them; the
team does not draw them.} Use at most one per round, after purchases are final
but before Phase 3, and always require the team to answer the card's question
aloud before continuing.

\begin{center}
\small
\begin{tabularx}{\textwidth}{@{}l X@{}}
\toprule
\textbf{Inject} & \textbf{Effect and when to play it} \\
\midrule
I1 Budget Cut & Withdraw two tokens; the team returns cards to that value. Play
on a team that finished early or bought without deliberating. \\
I2 Vendor End of Life & Name one purchased card as withdrawn. They may
substitute a Compensating Control if they can justify one. Play on a team
over-relying on a single product. \\
I3 Staff Departure & Suspend any purchase requiring ongoing human attention
(15, 13, 7, 9) until they name a successor. Play on a team that bought
monitoring with nobody to watch it. \\
I4 Restricted Windfall & Three extra tokens, must be spent this round on new
cards, unspent points lost. Play on a team that was too conservative. \\
\bottomrule
\end{tabularx}
\end{center}

Injects are for the second half of a course, once the base loop is internalized.
\textbf{Do not use them during any session you are collecting research data
from} --- they introduce variance that is not held constant across teams.

% =====================================================================
\section{Debrief}

Five to ten minutes, and the highest-value part of the session. Each scenario in
Appendix~\ref{app:scenarios} carries its own questions, which are sharper than
these because they name the specific trade-off that scenario was built around.
Use those first, then any of the general ones below. Ask three or four, not all:

\begin{itemize}
\item Which purchase did your team argue about most, and how did you settle it?
\item What did you want to buy and could not afford? What would you tell the
client about that gap?
\item Did anything you bought only help \emph{after} the attack succeeded? Was
that worth the points?
\item Two teams here bought different portfolios for similar threats. Why?
\item If your budget doubled, what would you add --- and why that, before
anything else?
\item Which card would you cut from this deck entirely?
\end{itemize}

\subsection{Misconceptions Worth Naming}

\begin{center}
\small
\begin{tabularx}{\textwidth}{@{}p{5.4cm} X@{}}
\toprule
\textbf{Common belief} & \textbf{What to say} \\
\midrule
Controls eliminate risk & They reduce likelihood or impact. Phase 3 exists to
show this; name it explicitly at debrief. \\
More expensive means better & Penetration testing costs three and is often the
wrong buy for a small organization. Cost tracks Implementation Group, not value. \\
Awareness training solves phishing & It is one layer. Email filtering (card 12)
is cheaper and does not depend on a tired person at 4pm. \\
Backups prevent breaches & They limit consequence. The data is still disclosed. \\
Everything must be fixed & Four risk treatments exist: avoid, mitigate,
transfer, accept. Cards 21, 26 and 27 exist for the other three. \\
Technical sophistication equals severity & Severity is who gets harmed. That is
why Phase 1 requires pointing at the Human Impact card. \\
\bottomrule
\end{tabularx}
\end{center}


% =====================================================================
\section{Closing the Session}

Discussion evaporates. A session that ends when the talking stops leaves
students with a pleasant memory and no artifact, and the translation step
below is the one clinic skill the rest of the game never asks for.

\subsection{The Last Three Minutes}

Ask every team for \textbf{one sentence: what would you actually tell this
organization?}

The constraint is that it must be in the client's words, not the deck's.
\say{Turn on multi-factor authentication} is a card name. \say{Nobody should be
able to move money on the strength of an email, even from the director} is a
recommendation. A student who can name a control but cannot say it to a nonprofit
director has not finished the work.

Go around the room. Write them on the board. They take ninety seconds and they
are the thing students remember in week eight.

\begin{note}{Why this and not a summary}
Summarizing is the facilitator's voice. This is theirs, and it is the only
moment in the activity where a team has to commit to a single position with no
cards to hide behind.
\end{note}

\subsection{Before You Leave the Room}

\begin{itemize}
\item \textbf{Export the session} if you are using the web version, or collect
the record sheets if you are on paper. Do this before anyone closes a laptop.
\item \textbf{Note which interventions you used}, at which rung. It is the only
way to tell later whether you are intervening too early.
\item \textbf{Reset or archive the board.} On paper this is shuffling the deck.
The web version has two resets, and picking the wrong one is annoying rather
than destructive. \emph{New round} clears threats, hands, placements and
purchases but leaves teams seated --- use it when the same class plays a second
scenario. \emph{End session and clear board} empties everything including team
membership, and keeps a timestamped copy first --- use it at the end of a class,
so the next one does not inherit a board that looks broken. A team can also undo
its own purchases mid-round with \emph{Clear our picks}, which leaves their hand
and threat intact.
\item \textbf{Do not score anything live.} Judging justifications while teams are
still in the room changes how they talk. Score from the export afterwards.
\end{itemize}

\subsection{What Comes Next in the Course}

The activity is an on-ramp, so it should hand off to something.

\textbf{Immediately after.} Students' one-sentence recommendations go in their own
notes. Ask them to keep them; they will be embarrassed by them later, which is
the point.

\textbf{Next session.} Two options, both cheap. Re-run the same scenario at a
different budget --- an \textsc{ig2} organization with ten points instead of six
--- and ask what changed. Or run a different scenario chosen to reward a control
the class ignored last time. If nobody bought logging in week three, run scenario
7, where detection is the only lever that works.

\textbf{Later in the term.} The one-sentence recommendation becomes a paragraph,
then a section of an actual client memo. When students write that memo, hand back
their sentence from week three. The distance between the two is the clearest
evidence of learning the course will produce, and it costs nothing to collect.

\textbf{In campaign play.} Do not reset. The same organization returns in a later
week with drift applied to any threat left unmitigated, and the record sheet is
what carries that forward. You are the bookkeeper across sessions
(Section~\ref{sec:campaign-manual}).

\subsection{If the Session Went Badly}

Two failure modes are worth naming, because both look like the activity not
working when they are actually diagnostic.

\textbf{Teams agreed too quickly and nobody argued.} The budget was too loose or
the deal too uniform. Cut to five points next time, or vary the piles between
teams so cross-team comparison has something to compare.

\textbf{Teams were overwhelmed and produced nothing.} Too many cards at once, or
you opened with V6. Drop to V3, deal six cards, pre-populate two cells of the
risk matrix, and come back to the full loop in a later week.

Neither is a reason to abandon the activity. Both are settings.

% =====================================================================
\section{Assessment}

If you are grading, score the \emph{justifications}, not the portfolio. Two
teams may correctly buy different things.

\begin{center}
\small
\begin{tabularx}{\textwidth}{@{}c l X@{}}
\toprule
\textbf{Score} & \textbf{Level} & \textbf{Description} \\
\midrule
0 & Assertion & Names a card with no reason. \say{We bought backups.} \\
1 & General reason & A reason that would apply to any threat. \say{Backups are
important for every organization.} \\
2 & Threat-specific & Links the card to this threat. \say{Ransomware encrypts
the donor database, so backups get it back.} \\
3 & Mechanism + safeguard & Links to this threat and names the mechanism or
safeguard. \say{Ransomware encrypts the donor database. Safeguard 11.4, an
isolated copy, means the attacker cannot reach the backup from the same
network.} \\
\bottomrule
\end{tabularx}
\end{center}

Expect mostly 1s in a first session and mostly 2s by the third. Level 3 requires
the card backs to have been read, which is itself a useful signal.

% =====================================================================
\section{Troubleshooting}

\begin{center}
\small
\begin{tabularx}{\textwidth}{@{}p{5cm} X@{}}
\toprule
\textbf{Problem} & \textbf{Fix} \\
\midrule
Teams finish far too fast & Deal fewer Foundational and more Extended cards, so
the budget binds harder. Or cut the budget to five. \\
Teams cannot finish in time & Deal six cards instead of eight. Pre-populate two
cells of the risk matrix. \\
Nobody uses the modifiers & Deal them face-up and say they are free. Most teams
do not notice they exist. \\
Every team buys the same three cards & Your deal is too uniform. Vary the piles
between teams so the cross-team comparison has something to compare. \\
Scenarios are implausible & Leave it. Plausibility is not the learning
objective, and auditing it costs time and morale. \\
One team is far ahead of the others & Play an inject (Section 9), or move them
to a 10-point budget and a harder scenario. \\
Discussion is polite and flat & The budget is too loose or the deal is too easy.
Six tokens against eight cards should hurt. \\
\bottomrule
\end{tabularx}
\end{center}

\newpage
% =====================================================================
\appendix
\section{Session Record Sheet}

\textit{Photocopy one per team per session. Retain across sessions for campaign
play --- you are the bookkeeper.}

\vspace{0.8em}
\noindent
\begin{tabularx}{\textwidth}{@{}p{3.4cm} X@{}}
\toprule
Team & \rule{0pt}{2.2em} \\
\midrule
Session / week & \rule{0pt}{2.2em} \\
\midrule
Organization & \rule{0pt}{2.2em} \\
\midrule
Budget issued & \rule{0pt}{2.2em} \\
\midrule
Drawn cards (4 dims) & \rule{0pt}{3em} \\
\midrule
Scenario (one line) & \rule{0pt}{3em} \\
\midrule
\textbf{Phase 1} placement & Likelihood: \rule{2.5cm}{0.4pt} \quad Severity:
\rule{2.5cm}{0.4pt}\\
\midrule
Severity justified by which Human Impact card? & \rule{0pt}{2.4em} \\
\midrule
\textbf{Phase 2} purchases and cost & \rule{0pt}{5em} \\
\midrule
Modifiers played / accepted? & \rule{0pt}{3em} \\
\midrule
Tokens spent & \rule{0pt}{2.2em} \\
\midrule
\textbf{Phase 3} placement & Likelihood: \rule{2.5cm}{0.4pt} \quad Severity:
\rule{2.5cm}{0.4pt}\\
\midrule
Stated residual risk & \rule{0pt}{3em} \\
\midrule
Justification score (0--3) & \rule{0pt}{2.2em} \\
\midrule
Deferred? Drift carried forward & \rule{0pt}{2.4em} \\
\midrule
Interventions used (rung) & \rule{0pt}{2.4em} \\
\midrule
Facilitator notes & \rule{0pt}{5em} \\
\bottomrule
\end{tabularx}

\newpage
% =====================================================================
%  scenarios.tex — Appendix C: Scenario Library
%  \input by facilitator_manual.tex
%
%  Ten ready-to-run scenarios. Each is engineered so that a DIFFERENT
%  control is the strongest answer. If every scenario rewarded the same
%  five cards, the game would teach one lesson repeatedly and students
%  would learn the answer rather than the reasoning.
%
%  All organizations are fictional. Never run a session against a real
%  client's systems.
% =====================================================================

\section{Scenario Library}
\label{app:scenarios}

Improvised scenarios are the most common reason a session goes flat, and
they also make cross-team comparison meaningless, since two teams arguing
about different situations are not disagreeing about anything. The ten
scenarios below are ready to read aloud.

Each carries a \textbf{teaching target}: the specific thing it is built to
surface. Each also carries an \textbf{expert key} naming which controls are
strong, partial, or weak \emph{for that situation}. The key is a scoring aid,
not an answer sheet --- a team that buys a ``partial'' card and defends it well
has done better work than one that buys a ``strong'' card without reasoning.

\begin{center}
\small
\begin{tabularx}{\textwidth}{@{}c X l c c@{}}
\toprule
\textbf{\#} & \textbf{Organization and teaching target} & \textbf{Budget} & \textbf{Tier} & \textbf{Care} \\
\midrule
1  & Food bank --- severity is consequence, not sophistication & 6 & Found. & \\
2  & Rural health clinic --- recovery beats prevention once it has happened & 6 & Found. & \\
3  & Domestic violence shelter --- when impact is physical safety & 6 & Found. & $\bullet$ \\
4  & Community land trust --- the boring control is the right one & 5 & Found. & \\
5  & School district --- you cannot buy your way out of a vendor's breach & 10 & Ext. & \\
6  & Municipal water utility --- the correct answer costs one point & 10 & Ext. & \\
7  & Legal aid office --- prevention fails, detection is the only lever & 10 & Ext. & $\bullet$ \\
8  & Public library --- you cannot protect what you do not know exists & 6 & Found. & \\
9  & Arts nonprofit --- what a grant paid for and nobody maintains & 6 & Camp. & \\
10 & Animal rescue --- the obligation that survives the incident & 10 & Ext. & \\
\bottomrule
\end{tabularx}
\end{center}

\vspace{0.4em}
{\small $\bullet$ Content that may land on a participant's personal experience.
Offer the redraw described in Section~7.4 \emph{before} reading it, addressed to
the team rather than to any individual.}

% --------- formatting macro ---------
\newcommand{\scen}[3]{%
  \vspace{1.0em}\noindent
  \textbf{#1. #2}\hfill{\small\textsc{#3}}\\[0.15em]
  \rule{\linewidth}{0.4pt}\\[0.25em]}
\newcommand{\sfld}[1]{\textit{#1}\hspace{0.4em}}

% =====================================================================
\scen{1}{Payroll Redirection at Harborview Food Bank}{Budget 6 \quad Foundational}
\sfld{Teaching target.} Severity is a property of who gets harmed, not of how
clever the attack was. The strongest control here is free and procedural.\\
\sfld{Read aloud.} Harborview Food Bank has three paid staff and about forty
volunteers. The bookkeeper works two days a week from home. Last Tuesday she
received an email that appeared to come from the board treasurer, in his usual
phrasing, asking her to update the payroll provider's bank details before
Friday's run. A follow-up voicemail in his voice confirmed it. Friday's payroll
is eighteen thousand dollars, roughly a month of operating cash.\\
\sfld{Suggested draw.} Method: Manipulation or Coercion $\cdot$ Impact: Financial
Wellbeing $\cdot$ Resources: Artificial Intelligence $\cdot$ Motivation:
Personal Gain\\
\sfld{Expert key.} \textbf{Strong:} Social Engineering Awareness (2), Deepfake
and AI Fraud Awareness (3), Email and Browser Protections (12). \textbf{Partial:}
Password and Authentication (5), Incident Response Planning (13).
\textbf{Weak here:} Encryption (14), Network Segmentation (18), Penetration
Testing (19) --- none of them touch a human being persuaded to type a new account
number.\\
\sfld{Debrief.}
\begin{itemize}
\item The attack was technically trivial and you rated it Very Severe. What made
it severe, and would it have been severe at a different organization?
\item Nobody can buy a callback rule. Which of your purchases would you trade
for one sentence in a written procedure?
\item The voice on the phone was cloned. Does awareness training still work when
familiarity stops being evidence?
\end{itemize}

% =====================================================================
\scen{2}{Ransomware at Cedar Ridge Health Clinic}{Budget 6 \quad Foundational}
\sfld{Teaching target.} Once an incident has happened, prevention cards are worth
nothing and recovery cards are worth everything. Teams that spent all six points
on prevention have to explain their Monday morning.\\
\sfld{Read aloud.} Cedar Ridge is a two-provider clinic serving about nine hundred
patients in a county with no other primary care. On Monday morning the front
desk cannot open the scheduling system. Every file on the shared drive has a new
extension, and a text file on the desktop asks for payment in cryptocurrency.
The practice manager thinks there is a backup, but it is on a drive plugged into
the same server, and nobody has ever restored from it.\\
\sfld{Suggested draw.} Method: Technological Attack $\cdot$ Impact: Physical
Wellbeing $\cdot$ Resources: Money $\cdot$ Motivation: Personal Gain\\
\sfld{Expert key.} \textbf{Strong:} Backup and Recovery (6) --- specifically the
isolated copy, since a backup on the same server is now encrypted too;
Incident Response Planning (13); Endpoint Detection and Response (16).
\textbf{Partial:} Software Updates (1), Email and Browser Protections (12),
Cyber Insurance (21). \textbf{Weak here:} Penetration Testing (19) --- an annual
report does not help on the morning it happens.\\
\sfld{Debrief.}
\begin{itemize}
\item The backup existed and was useless. What is the difference between having
a backup and having a recovery?
\item Impact is Physical Wellbeing, not Financial. What does a clinic without
scheduling actually do to a patient?
\item If you bought only prevention, describe your Monday.
\end{itemize}

% =====================================================================
\scen{3}{A Laptop Leaves the Building --- Wilder House}{Budget 6 \quad Foundational}
\sfld{Teaching target.} When Human Impact is physical safety, ordinary controls
carry stakes students have not previously attached to them.\\
\sfld{Handle with care.} This scenario involves a shelter and the safety of people
escaping violence. Offer the redraw to the team before reading it, and do not
require anyone to narrate.\\
\sfld{Read aloud.} Wilder House is a twelve-bed shelter. A caseworker's laptop was
taken from her car in a grocery store parking lot on Saturday. It holds an
offline copy of the intake spreadsheet: names, dates of birth, the schools
children attend, and current addresses for eleven families. The laptop has a
login password. Nothing else.\\
\sfld{Suggested draw.} Method: Physical Attack $\cdot$ Impact: Physical Wellbeing
$\cdot$ Resources: Entitlement $\cdot$ Motivation: Curiosity\\
\sfld{Expert key.} \textbf{Strong:} Encryption (14) --- full-disk encryption makes
this a hardware loss rather than a disclosure; Physical Security (8); Identity
and Access Management (4), on why an offline copy existed at all.
\textbf{Partial:} Asset Inventory (9), Incident Response Planning (13).
\textbf{Weak here:} Audit Logging (15), Network Segmentation (18) --- the device
is gone and off the network.\\
\sfld{Debrief.}
\begin{itemize}
\item A login password and full-disk encryption feel similar to a non-specialist.
What is the actual difference to whoever has the laptop now?
\item The spreadsheet was a convenience copy. Which control stops a convenience
copy from existing, and what does the caseworker lose when it does?
\item Who has to be told, how quickly, and what do you say to eleven families?
\end{itemize}

% =====================================================================
\scen{4}{The Volunteer Who Left in 2021 --- Fairmount Land Trust}{Budget 5 \quad Foundational}
\sfld{Teaching target.} The strongest answer is the least interesting card in the
deck. Teams reliably overlook it and buy something exciting instead.\\
\sfld{Read aloud.} Fairmount Land Trust has two staff and a rotating pool of
volunteers. A volunteer who managed the donation platform in 2021 moved out of
state and stopped responding. Last month a two-hundred-dollar test transaction
appeared, then reversed. The executive director does not know how many people
still have logins, and the platform bills annually to a credit card nobody can
locate the statement for.\\
\sfld{Suggested draw.} Method: Processes $\cdot$ Impact:
Financial Wellbeing $\cdot$ Resources: Entitlement $\cdot$ Motivation: Personal
Gain\\
\sfld{Expert key.} \textbf{Strong:} Account Lifecycle Management (7), Identity and
Access Management (4). \textbf{Partial:} Asset Inventory (9), Software Inventory
(10), Audit Logging (15). \textbf{Weak here:} Social Engineering Awareness (2)
--- nobody was tricked; a door was left open.\\
\sfld{Debrief.}
\begin{itemize}
\item Nobody attacked this organization. Is this a security incident?
\item The right answer costs one point and takes an afternoon. Why is it the one
nobody does?
\item How would this organization even produce a list of who has access?
\end{itemize}

% =====================================================================
\scen{5}{The Vendor's Breach --- Northbridge School District}{Budget 10 \quad Extended}
\sfld{Teaching target.} Sometimes nothing you can buy internally helps, and the
only real control was written a year earlier in a contract.\\
\sfld{Read aloud.} Northbridge is a rural district of four schools. Their bus
routing and student scheduling run on a hosted platform used by two hundred
districts. On Thursday a reporter calls the superintendent for comment on a
breach at that vendor. The district has had no notification. The vendor's
support line has a recorded message. The data includes student names, home
addresses, and bus stop times.\\
\sfld{Suggested draw.} Method: Indirect Attack $\cdot$ Impact: Relationships
$\cdot$ Resources: Technical Expertise $\cdot$ Motivation: Personal Gain\\
\sfld{Expert key.} \textbf{Strong:} Third-Party and Contractual Controls (17),
Incident Response Planning (13). \textbf{Partial:} Asset Inventory (9) extended
to data location, Cyber Insurance (21). \textbf{Weak here:} nearly everything
else --- Endpoint Detection, Segmentation, Secure Configuration and Software
Updates all protect systems the district does not own.\\
\sfld{Debrief.}
\begin{itemize}
\item Half your deck is useless here. What does that tell you about where a small
organization's real risk lives?
\item The superintendent learned from a reporter. What clause would have changed
that, and when would it have had to be written?
\item Bus stop times plus home addresses. Which Human Impact card is this really?
\end{itemize}

% =====================================================================
\scen{6}{Debug Mode in Production --- Millbrook Water Authority}{Budget 10 \quad Extended}
\sfld{Teaching target.} A one-point Foundational card beats every expensive option
on the table. Teams with ten points often overspend here.\\
\sfld{Read aloud.} Millbrook Water Authority serves eleven thousand residents.
Their public site posts water quality reports and lets residents pay bills. A
contractor built it four years ago and moved on. A resident emails to say that
adding a parameter to a URL shows an administrative page listing every account
and a button labelled ``post notice.'' The page has no login.\\
\sfld{Suggested draw.} Method: Processes $\cdot$ Impact:
Relationships $\cdot$ Resources: Technical Expertise $\cdot$ Motivation: Politics\\
\sfld{Expert key.} \textbf{Strong:} Secure Configuration (11) --- debug mode in
production is a configuration failure; Identity and Access Management (4).
\textbf{Partial:} Penetration Testing (19) would have found it, at three points;
Software Inventory (10); Third-Party and Contractual (17), on the contractor who
left. \textbf{Weak here:} Social Engineering Awareness (2), Backup and Recovery
(6).\\
\sfld{Debrief.}
\begin{itemize}
\item You had ten points. What is the cheapest purchase that actually closes this?
\item Penetration testing would have found it. Was it worth three points here, and
what would have to be true first?
\item A utility can post public notices. What is the worst version of this that
does not involve stealing anything?
\end{itemize}

% =====================================================================
\scen{7}{The Trusted Paralegal --- Eastgate Legal Aid}{Budget 10 \quad Extended}
\sfld{Teaching target.} Every preventive control in the deck fails, because the
person had legitimate access. Detection is the only remaining lever.\\
\sfld{Handle with care.} Involves surveillance of a person by someone they trusted.
Offer the redraw.\\
\sfld{Read aloud.} Eastgate Legal Aid handles housing and immigration matters. A
paralegal of six years has legitimate access to the case system. A client
reports that her landlord seems to know the contents of a filing that has not
been served. The paralegal and the landlord attend the same church. There is no
evidence of an intrusion, and nothing in the system is broken.\\
\sfld{Suggested draw.} Method: Processes $\cdot$ Impact:
Emotional Wellbeing $\cdot$ Resources: Entitlement $\cdot$ Motivation: Desire or
Obsession\\
\sfld{Expert key.} \textbf{Strong:} Audit Logging and SIEM (15) --- who read which
file and when is the only question that can be answered; Identity and Access
Management (4) with Least Privilege (24); Managed Detection and Response (20) if
affordable. \textbf{Partial:} Incident Response Planning (13), Third-Party and
Contractual (17) on employment terms. \textbf{Weak here:} Email and Browser
Protections (12), Encryption (14), Software Updates (1) --- the access was
authorized.\\
\sfld{Debrief.}
\begin{itemize}
\item Which of your purchases would have stopped this? If none, what would you
tell the client you can offer instead?
\item Logging tells you afterward. Is a control that only works afterward worth
two of ten points?
\item Least Privilege has a cost that is not measured in points. What does the
paralegal lose, and is the trade right?
\end{itemize}

% =====================================================================
\scen{8}{The Camera Nobody Remembered --- Delridge Public Library}{Budget 6 \quad Foundational}
\sfld{Teaching target.} Asset inventory is a precondition for everything else.
Teams cannot defend what they have not enumerated.\\
\sfld{Read aloud.} Delridge Public Library has public wifi, twenty patron
computers, a self-checkout system, and four security cameras installed by a
since-dissolved vendor in 2016. The county IT contractor reports that traffic is
leaving the library network at three in the morning, and traces it to a device
nobody on staff can identify. The library's network is flat: everything is on
the same wifi, including the circulation desk.\\
\sfld{Suggested draw.} Method: Technological Attack $\cdot$ Impact: Relationships
$\cdot$ Resources: Technical Expertise $\cdot$ Motivation: Curiosity\\
\sfld{Expert key.} \textbf{Strong:} Asset Inventory (9), Network Segmentation (18)
if the budget reaches it, Secure Configuration (11) for the 2016 default
credentials. \textbf{Partial:} Software Inventory (10), Audit Logging (15).
\textbf{Weak here:} Social Engineering Awareness (2), Backup and Recovery (6).\\
\sfld{Debrief.}
\begin{itemize}
\item At six points, segmentation may be out of reach. What do you buy first, and
what do you tell the library board about the rest?
\item The camera has been there for a decade. Which control would have caught it
in year one, and how much does that control cost?
\item Patron borrowing records are legally protected in many states. Does that
change your ranking?
\end{itemize}

% =====================================================================
\scen{9}{What the Grant Paid For --- Ridgeway Arts Collective}{Budget 6 \quad Campaign}
\sfld{Teaching target.} Built for multi-engagement play. The right move is often to
defer, and the card that lets you defer is the one that makes the cost of
deferring visible.\\
\sfld{Read aloud.} Ridgeway Arts received a technology grant in 2019 that paid for
a donor database, a website, and a year of support. The grant ended. The database
runs a version the vendor stopped patching in 2022. Migrating costs about four
thousand dollars the collective does not have. The board meets quarterly. Nothing
has gone wrong yet.\\
\sfld{Suggested draw.} Method: Technological Attack $\cdot$ Impact: Financial
Wellbeing $\cdot$ Resources: Technical Expertise $\cdot$ Motivation: Personal
Gain\\
\sfld{Expert key.} \textbf{Strong:} Software Inventory (10), Deferred Investment
(28) played honestly with a stated exposure, Backup and Recovery (6).
\textbf{Partial:} Network Segmentation (18) to isolate the unsupported system,
Compensating Control (27), Risk Acceptance (26) with a named owner.
\textbf{Weak here:} Software Updates (1) --- there are no patches to apply, which
is the whole problem, and teams reliably buy it anyway.\\
\sfld{Debrief.}
\begin{itemize}
\item Someone bought Software Updates. What happens when the vendor has stopped
shipping them?
\item If you deferred, name the exposure and say for how long. Would you put that
in writing to the board?
\item Revisit this in a later session and drift the threat. Was deferring still
the right call?
\end{itemize}

% =====================================================================
\scen{10}{Donor Records in a Shared Mailbox --- Second Chance Animal Rescue}{Budget 10 \quad Extended}
\sfld{Teaching target.} Obligations survive the incident. Notification and
disclosure are not technical controls and cannot be bought after the fact.\\
\sfld{Read aloud.} Second Chance runs on donations and a shared mailbox that six
volunteers access with the same password. On Friday the mailbox began sending
donation appeals to the entire contact list from an address that is not theirs.
The mailbox holds four years of correspondence, including scanned cheques with
routing numbers and a spreadsheet of major donors.\\
\sfld{Suggested draw.} Method: Manipulation or Coercion $\cdot$ Impact: Financial
Wellbeing $\cdot$ Resources: Money $\cdot$ Motivation: Personal Gain\\
\sfld{Expert key.} \textbf{Strong:} Password and Authentication (5) --- shared
credentials plus no multi-factor is the whole story; Account Lifecycle
Management (7); Incident Response Planning (13) for the notification duty.
\textbf{Partial:} Email and Browser Protections (12), Audit Logging (15), Cyber
Insurance (21). \textbf{Weak here:} Network Segmentation (18), Endpoint Detection
(16) --- nothing was installed on any machine.\\
\sfld{Debrief.}
\begin{itemize}
\item Six people share one password because it is genuinely convenient. What do
you replace it with that they will actually use? (This is the Human Factors card
in disguise.)
\item Which of your purchases helps you tell four thousand donors what happened?
\item Insurance may cover the cost. Does it cover the donor who stops giving?
\end{itemize}

% =====================================================================
\vspace{1.2em}
\subsection*{Writing Your Own}

A scenario that works has four properties. It names an organization concrete
enough to picture. It states a consequence in human terms rather than technical
ones. It has \emph{more than one} defensible answer, or teams agree in ninety
seconds. And it has at least one card that looks obviously right and is not ---
that card is where the argument happens.

Test a draft by asking what a team would buy with it, then asking whether a
different team could reasonably buy something else. If the answer to the second
question is no, the scenario is a quiz question rather than a game.


\newpage
\section{Facilitator Self-Check}

Run this before your first session and after your third.

\vspace{0.6em}
\noindent\textbf{Before the first session}
\begin{itemize}
\item[$\square$] I have read Section 3 (Quick Reference) and Section 7.
\item[$\square$] Decks are sorted into Foundational / Extended / Advanced /
modifiers.
\item[$\square$] I have physical tokens, not a tally sheet.
\item[$\square$] I know the four immediate-intervention situations in 7.4.
\item[$\square$] I have a redraw available for participants who want a different
Impact card, and I will offer it to the team rather than to an individual.
\item[$\square$] I am prepared to explain Phase 1 only, then start.
\end{itemize}

\vspace{0.6em}
\noindent\textbf{After the third session}
\begin{itemize}
\item[$\square$] Did teams argue? If not, tighten the budget or vary the deal.
\item[$\square$] Did I intervene above Rung 2 more than twice per table? If so, I
am intervening too early.
\item[$\square$] Did any team move a threat to Not Likely / Not Severe unchallenged?
\item[$\square$] Did non-majors speak as much as majors?
\item[$\square$] Are justification scores rising across sessions?
\item[$\square$] Did I supply an answer that a peer would have supplied within
five seconds?
\end{itemize}

\vspace{1.4em}
\begin{note}{Feedback}
This manual is a working document. Where a rule did not survive contact with
your classroom, change it and note why --- re-tiering costs, adjusting the
budget, and cutting cards are all expected local adaptations.
\end{note}

\end{document}
